\begin{inplace}{figure}
\centering\includegraphics[width=\linewidth]{crystal-examples.pdf}
\caption{The first four distinct generated formulas, before and after CHGNet relaxation, in a common orientation and at the same physical scale. Saturated spheres mark sites in the central cell and pale spheres mark periodic images. Thin rods outline the cell. Sticks mark a pair when its distance is within 1.18 times the nearest periodic distance of either endpoint. They indicate geometric proximity. The LiPd$_3$ example reaches the optimizer step limit with a 0.313\,\AA\ Pd--Pd contact, indicated on the relaxed panel.}
\label{fig:structures}
\end{inplace}

We relaxed the cell and all sites at zero external pressure using the published CHGNet 0.3.0 weights \cite{chgnet}, ASE FIRE, a Fr\'echet cell filter, a generalized force threshold of 0.1\,eV/\AA\ and a cap of 100 steps. Controls included 12 randomly selected held-out structures and one available same-reduced-formula MP-20 reference for each of 23 draws. Reference selection prioritized the test split, then validation and training, with site count and material ID breaking ties. Each reference underwent the same relaxation. The surrogate was itself trained on Materials Project trajectories. This screen therefore provides model-dependent evidence.

Across the 32 generated structures, median maximum atomic force fell from 32.44 to 0.097\,eV/\AA, and median nearest contact increased from 1.44 to 2.50\,\AA. Median energy change was $-4.42$\,eV/atom. Twenty-three runs met the optimizer's combined cell-and-position stopping rule. Seventeen ended with maximum physical atomic force at or below 0.1\,eV/\AA, as did all 12 random controls. A generated cell took a median 0.039 seconds for the initial evaluation and 3.12 seconds in the optimizer on CPU. Median recorded wall time was 3.16 seconds. Prospective ranking was absent from both the proposal timer and this consecutive-draw screen. Seven retained a contact below 60\% of the sum of the pair's covalent radii, and three retained a contact below 0.5\,\AA. In Fig.~\ref{fig:structures}, MgSm$_2$ relaxed to within 0.004\,eV/atom of its selected reference, while the LiPd$_3$ cell retained a severe overlap.

\begin{inplace}{figure}
\centering\includegraphics[width=.93\linewidth]{revision-relaxation.pdf}
\caption{CHGNet relaxation of 32 consecutive proposals. Each marker in panels a and b pairs one initial and final geometry. Dashed diagonals indicate no change. Dotted lines mark $0.1$\,eV/\AA\ and a contact-to-radii ratio of $0.60$. Panel c compares each of the 23 available same-formula references after identical relaxation. Positive energies favor the reference.}
\label{fig:screen}
\end{inplace}


Of the 23 candidates with a same-formula reference, 20 ended above it in predicted energy. The median excess was 0.041\,eV/atom, with a range from $-0.103$ to $+9.967$\,eV/atom (Fig.~\ref{fig:screen}). We also relaxed one deposited elemental reference per constituent element, 44 in total. We chose each by lowest deposited energy above hull, then formation energy, site count and ID. All 44 converged. Comparing each generated structure with its composition-weighted elemental references gave 19/32 below the mixture and 13/32 above, with a median difference of $-0.046$\,eV/atom. This comparison omits competing binary phases. A converged geometry or negative energy relative to elemental references alone cannot establish survival against all decomposition pathways. Bulk stability, dynamical stability and synthesis remain unverified.

We also compared all 1,024 unrelaxed proposals using species-aware \texttt{StructureMatcher.fit} at the same lattice and angle tolerances as the conditional task. Counting the first representative of each matching group in draw order produced 1,014 representatives. Searching the complete MP-20 release and all 256,259 eligible WBM cells left 940 proposals without a match in either collection. Of the original 1,024 cells, 859 were simultaneously a greedy representative, unmatched in both databases, and free of contacts shorter than 0.7\,\AA. LeMat-GenBench additionally assesses validity and stability against a self-consistent competing-phase hull \cite{lemat}. Database nonmatches alone do not establish stability or discovery.
